\section{Open items}
\label{sec:open}

The comparisons leave seven open items, each stated with the measurement
or calculation that would settle it. Five are summarized here with
pointers to their sections; the high-redshift kinematics and the Milky Way
mass follow in full.

\paragraph{Rotation-curve normalization.} On SPARC the stock law trails
MOND with a fitted $a_0$ by 108 to 131 in $\chi^2$ and MOND at
$1.2\times10^{-10}\,\mathrm{m\,s^{-2}}$ by 26 to 34. The law shape is
shared, so the gap measures the normalization. The fitted value spreads
by 0.25~dex between galaxy classes and contains the cosmological $a_0$
(Section~\ref{sec:sparcthree}). A 2\% determination of $a_0$ from galaxies
with absolute distances, with correlated distance and stellar-population
errors, tests the normalization and the Hubble-constant relation of
Section~\ref{sec:h0galaxies}.

\paragraph{Isolated-galaxy lensing.} MOND fits the KiDS isolated-lens
profiles with $\chi^2=110$ and no halo parameter. The support-limited stock
reaches 132 only at a third of the standard concentration
(Section~\ref{sec:kidsthree}). The hot-gas mass of the lenses and the
profile of the uncarried stock outside the support decide this comparison.
A forward model of the stacked profiles with the support radius measured
from the circumgalactic gas, and the same environment term for all models,
is the next test.

\paragraph{Antlia 2.} The best stock remnant has $3.48\,\kms$ against the
measured $6.0\,\kms$ (Section~\ref{sec:antlia}). MOND with the external
field also falls short. A mock observation with the actual membership,
gradient, and selection procedure is required.

\paragraph{Release normalization.} The archived reference gate of
$17\,\kms$ gave a release integral about 20\% below the likelihood target,
and the likelihood selects $14\,\kms$. The cooling criterion and the
no-feedback run, computed afterwards, give larger vacuum fractions. The
feedback-resolved occupation was adopted from published simulations after
the target and these gates were known, and it lands on the target with
none of its inputs adjusted. A retention fraction measured in advance
from converged runs with feedback and photoheating, under a label prescription fixed beforehand, is the
decisive form of this test.

\paragraph{Cluster outliers.} A1795 and A2029 require transport of
baryons out of the lensing aperture, and A644 and ZW1215 retain
reconstruction uncertainty (Section~\ref{sec:clusters}).

\subsection{High-redshift kinematics and the direction of a relaxation correction}
The reference line projection agrees with the published coefficients
(Section~\ref{sec:redshift}). The catalogue residuals retain a trend that
the $H(z)$ scaling does not remove, and a constant $a_0$ leaves a slightly
smaller one. The catalogue diagnostics use the exponential-allocation
normalization $a_0=1.2588\times10^{-10}\,\mathrm{m\,s^{-2}}$.

The common cosmology was propagated to 73 KMOS3D, 21 KGES, and 169 KROSS
objects at three radii. At the KMOS3D outer radius, the log-acceleration
residual slope is $-0.245$ dex per unit redshift with constant $a_0$ and
$-0.313$ with the common $H(z)$ law. The latter has median log-velocity
residual $-0.087$ dex. The reference $z=2.3$ example gives acceleration
boosts about 2.02 versus 1.39 for the compared evolving and fixed scales.
These are catalogue central-value diagnostics with specified baryon force
models; they do not constitute raw-cube likelihoods. The 112-point
reference projection remains distinct from a regression on the actual
79-object MUSE selection. The common cosmology has not removed the
high-redshift residual trend. \source{U1, S01}

Passive relaxation cannot automatically supply a reduction. If
$\dot n=\Gamma(q-n)$ with $\Gamma>0$, $\dot q\le0$, and initially
$n\ge q$, then $n$ remains above the target. At fixed baryons a lag behind
declining $H$ therefore has the wrong sign to suppress an initially
equilibrated response. Underoccupied assembly states require their own
history. Relabeling an unchanged stock leaves its total stress unchanged;
the tested adiabatic cosmological mass loss gives an inferred scale ratio
0.9183 from $z=2.3$, compared with $H_0/H(2.3)=0.2972$.
These constraints identify the contact and preparation history as physical
requirements. \source{D1}

For the selected two-pole quantum response, the noise enhancement
$\mathcal F$ and mean delay $\bar\tau$ obey
$\mathcal F H\bar\tau\ge4$. The illustrative three-percent tracking
condition at occupation argument $x=1$ gives $\mathcal F\gtrsim158$,
requiring about $6.47\times10^4$ independent cells for three-percent
force rms. Actual windows, correlations, and gravitational coupling
remain unspecified, so this is a contact-model constraint.
\source{D1; T1, N.8}

\subsection{Milky Way mass and a joint satellite fit}
\label{sec:jointtest}
The available
moving-LMC calculations use a Galactic mass within 50 kpc about $17\%$
above the Brooks et al. inference in the stock model, and about $36\%$
above it in the NFW model. A joint Milky Way--LMC fit would alter the
satellite pericentres and the stripping histories. Those changes should
be propagated to both Crater II and Antlia 2 before assigning a combined
evidence ratio. \source{G4}\cite{brooks}

\subsection{A decisive program}
The open items define a program:
\begin{enumerate}
\item Infer the Milky Way and LMC jointly using streams, halo reflex
kinematics, and a common source-conserving force model. Predict additional
stream observables from the resulting mass distributions.
\item Evolve the dwarf populations with resolved self-gravity and
prespecified progenitor priors. Apply the actual stellar selection,
velocity-gradient, and membership inference to their mock observations.
Include concentration scatter and physically motivated cores in the
dark-matter comparison.
\item Fit the intermediate-redshift galaxy data directly with the
cosmological $H(z)$ from the same model. Preserve correlations between
baryonic masses, pressure support, geometry, and acceleration.
\item Measure retention and source-label loss in converged hydrodynamic
runs with supernova feedback, photoheating and a label prescription fixed
in advance, and read the retention gate from gas retention. Use that release history
with the existing positive-allocation perturbation implementation and
a matched full cosmological likelihood.
\item Forward-model the isolated-lens profiles with the support radius
measured from the circumgalactic gas, and compare MOND, halos, and stock
with one environment term.
\end{enumerate}
Each step tests an existing physical commitment of the construction.

